\documentclass[11pt,a4paper]{ltjsarticle}
\usepackage{amsmath,amssymb}
\usepackage{graphicx}
\usepackage{booktabs}
\usepackage{array}
\usepackage{luatexja-fontspec}
\setmainjfont{HaranoAjiMincho-Regular}
\setsansjfont{HaranoAjiGothic-Medium}
\linespread{1.05}\selectfont

% Recommended PDF filename:
% 2611193_HeizoNakai_visual_media_processing1_dai4kai.pdf
% Example build command:
% latexmk -pdfdvi -jobname=2611193_HeizoNakai_visual_media_processing1_dai4kai main.tex

\begin{document}

\begin{center}
{\LARGE \textbf{視覚メディア処理I 第4回レポート}}
\end{center}
\vspace{1em}

\begin{flushright}
提出日: \today \\
学籍番号: 2611193 \\
氏名: 中井 平蔵
\end{flushright}

JPEGは、人間の視覚特性を利用して、見た目への影響が小さい情報を減らす圧縮方式である。
処理の流れとしては、まずRGB画像をYCbCr色空間に変換し、色差成分のデータ量を削減する。
その後、画像を8×8ピクセルのブロックに分割し、各ブロックに対してDCTと量子化を行う。
最後に、ジグザグスキャン、ランレングス符号化、ハフマン符号化によってさらにデータ量を削減する。
JPEG圧縮方式の中で、特に優れている技術として、私は「DCTと量子化」と「クロマサブサンプリング」の2つを選ぶ。

\subsection*{1. DCTと量子化}

DCTは、画像を8×8ピクセルのブロックに分割し、
そのブロック内の輝度値の分布を周波数成分として表現する技術である。
画像の周波数成分とは、輝度値の空間的な変化の細かさを表す。
空や海のように輝度値がなめらかに変化する領域は低周波成分を多く含み、
文字、細かい模様のように輝度値が急激に変化する領域は高周波成分を多く含む。
DCTを行うと、8×8ブロック内の輝度値の分布を、各周波数成分の強さを表すDCT係数として表すことができる。

JPEGでは、このDCT係数に対して量子化を行う。
量子化とは、DCT係数をある値で割り、細かい値を丸める処理である。
人間の視覚では、画像全体の印象を左右する低周波成分の変化は知覚されやすい一方、
細かい模様や微細な輝度変化を表す高周波成分の変化は比較的知覚されにくい。
そのためJPEGでは、低周波成分は小さい値で割ることで情報を細かく残し、
高周波成分は大きい値で割ることで強く削減する。
これにより、視覚的に重要な低周波成分を保ちながら、人間が変化を知覚しにくい高周波成分を効率的に削減できる。
結果として、JPEGは見た目への影響を抑えつつ、画像のデータ量を大幅に減らすことができる。

この技術が優れている理由は、人間の視覚特性をうまく利用している点にある。
自然画像では、隣り合うピクセルの値が大きく変化しないことが多く、低周波成分を多く含む。
DCTを使うことで、この重要な低周波成分と、削減しやすい高周波成分を分離できる。
そして量子化によって、高周波成分を積極的に削減することで、圧縮効率を高めることができる。

効果が明確に現れる画像の例としては、空、海を含む風景写真が挙げられる。
これらの画像はなめらかな領域が多いため、DCT係数が低周波成分に集中しやすい。
そのため、高周波成分を量子化によって削減しても見た目の劣化が比較的小さく、JPEG圧縮が効果的に働く。
一方で、圧縮率を高くしすぎると、8×8ブロックの境界が目立つブロックノイズや、
輪郭の周りに生じるリンギングノイズが現れる。

\subsection*{2. クロマサブサンプリング}

クロマサブサンプリングは、明るさとは別の「色のずれ」の情報を少なくすることでデータ量を削減する技術である。
JPEGでは、RGB画像をYCbCr色空間に変換する。
Yは輝度成分を表し、画像の明るさの情報を担う。
一方、CbとCrは色差成分であり、輝度成分Yから分離された色の情報を表す。
Cbは主に青方向の色差を、Crは主に赤方向の色差を表し、画像の色味を構成する役割を持つ。

人間の視覚では、輝度値の変化は知覚されやすい一方、色の細かな変化は比較的知覚されにくい。
そのため、色味の情報の解像度を多少下げても、人間は違いを知覚しにくい場合が多い。
JPEGではこの性質を利用し、輝度成分は高い解像度で保持しながら、色の情報の解像度だけを下げる。

代表的な方式に4:2:0や4:2:2がある。
これらは、輝度成分Yに対して、色差成分CbとCrをどの程度の密度で記録するかを表す方式である。
4:2:2では、横方向に隣り合う画素で同じ色差成分を共有することで、色差成分のデータ量を減らす。
一方、4:2:0では、横方向と縦方向の両方で色成分を共有するため、
4:2:2よりもさらにデータ量を削減できる。
そのため、色差成分のデータ量が少なくなる分、
色の境界部分では色が周囲に広がったように見える色のにじみが生じやすくなる。
しかし、どちらの方式でも輝度成分は残されるため、画像全体の形や輪郭は比較的保たれる。

この技術が優れている理由は、人間の視覚特性を利用して、
画質への影響が小さい部分を選択的に削減している点である。
単純に画像全体の解像度を下げると、明るさや輪郭の情報まで失われてしまい、画像がぼやけて見える。
しかし、クロマサブサンプリングでは、輝度情報を保ち、色の情報だけを削減する。
そのため、見た目の画質を保ちながら、効率よくデータ量を減らすことができる。

効果が明確に現れる画像としては、風景写真、人物写真、食べ物の写真などの自然画像が挙げられる。
これらの画像では、色がなめらかに変化する領域が多いため、色差成分の記録密度を下げても違和感が生じにくい。
一方で、白地に赤い文字が書かれた画像や、青と赤の細い線が隣り合う画像では、
文字や線の周辺に色のにじみが生じることがある。

\end{document}
